\subsection{Section summary}

This section defined the data foundation of the thesis. This thesis works on a 977-paper ingested corpus of histopathology-focused, English papers downloaded from PMC Open Access dataset. Two subsets of papers are derived from this ingested corpus: a 27-paper subset for document-extraction evaluation (\pinktodo{Section 3.5, document-extraction}, and another 15-paper subset selected by an ILP script based on paper relatedness. \pinktodo{Change the number here: 273} source cases were selected from the 15-paper subset, \pinktodo{change this number too: 1\,596} silver findings were generated by Claude Opus 4.7, then these findings were used for knowledge-extraction threshold calibration. The two paper subsets are completely disjoint: they do not share any papers so that improvements in the measurements of one subset does not affect the evaluation of the other subsystem.

This chapter also clarified the annotation status of the knowledge-extraction reference labels. These labels were LLM-generated, not expert-annotated, for scalability. These LLM-generated silver-labels, therefore, were not treated as ground truth. 

Overall, this chapter established what data the thesis uses, how the calibration and evaluation subsets were generated from the initial data, and the limitations of the data design. The following sections and chapters build on this foundation: \pinktodo{Section 3.3, Document extraction} describes document extraction, \pinktodo{Section 3.4, Knowledge extraction} describes knowledge extraction from the extracted text from the document-extraction pipeline, \pinktodo{Section 3.6, Evaluation Methodology?} defines evaluation metrics, \pinktodo{Chapter 4, Results} reports the results, and \pinktodo{Chapter 5, Discussion} discusses the implications of data limitations on the interpretation of the results.